% META: {"id": "1NSI-RES-ME-003", "chapitre": "1NSI-RESEAUX", "type_objet": "methode", "capacites": ["P-ARCH-02C"], "capacites_codes": ["C3"], "methodes": ["M3"], "mode_creation": "ex_nihilo", "sources_inspiration": [], "status": "needs_review"}
% BEGIN-VERIFY
% def liens_non_reciproques(reseau):
%     defauts = []
%     for machine, voisins in reseau.items():
%         for voisin in voisins:
%             if machine not in reseau.get(voisin, []):
%                 defauts.append((machine, voisin))
%     return defauts
% def peuvent_communiquer(reseau, depart, arrivee, visites=None):
%     if visites is None:
%         visites = set()
%     if depart == arrivee:
%         return True
%     visites.add(depart)
%     for voisin in reseau.get(depart, []):
%         if voisin not in visites:
%             if peuvent_communiquer(reseau, voisin, arrivee, visites):
%                 return True
%     return False
% reseau = {
%     "PC": ["S1"],
%     "S1": ["PC", "S2", "S3"],
%     "S2": ["S1", "S3"],
%     "S3": ["S1", "S2", "Serveur"],
%     "Serveur": ["S3"],
% }
% assert liens_non_reciproques(reseau) == []
% visites = set()
% assert peuvent_communiquer(reseau, "PC", "Serveur", visites) is True
% assert "S1" in visites
% brouillon = {"Box": ["PC", "TV"], "PC": ["Box"], "TV": []}
% assert liens_non_reciproques(brouillon) == [("Box", "TV")]
% assert peuvent_communiquer(brouillon, "Box", "TV") is True
% assert peuvent_communiquer(brouillon, "TV", "Box") is False
% END-VERIFY
\begin{fichemethode}{M3}{Modéliser un réseau en dictionnaire, puis chercher un chemin en marquant les visités}

\textbf{Quand l'utiliser ?} Quand l'énoncé décrit un réseau (machines, commutateurs,
câbles) et demande de le représenter en Python, ou de dire si deux machines peuvent
communiquer.

\textbf{Pas à pas --- construire le dictionnaire :}
\begin{enumerate}
  \item \textbf{Une clé par machine}, y compris les commutateurs et les routeurs : ce sont
        des sommets comme les autres. Une machine sans aucune connexion a une liste vide,
        elle n'est pas oubliée.
  \item \textbf{Chaque câble apparaît deux fois} : si \lstinline{"A"} figure dans la liste de
        \lstinline{"B"}, alors \lstinline{"B"} figure dans celle de \lstinline{"A"}. Vérifie
        cette réciprocité avant toute recherche ; un lien noté d'un seul côté donne des
        réponses différentes selon le sens de la question.
\end{enumerate}

\textbf{Pas à pas --- chercher un chemin à la main :}
\begin{enumerate}
  \item \textbf{Tiens à jour l'ensemble des visités.} Écris-le et complète-le à chaque
        machine explorée : c'est lui qui empêche de tourner en rond dans un cycle.
  \item \textbf{Depuis la machine courante, explore les voisins non visités}, un par un.
        Dès que l'arrivée est atteinte, la réponse est \lstinline{True}.
  \item \textbf{Si tous les voisins ont été explorés sans succès}, la réponse est
        \lstinline{False} : il n'y a pas de chemin. Et la recherche s'arrête toujours, parce
        que chaque machine n'est visitée qu'une fois.
\end{enumerate}

\textbf{Exemple rédigé.} Un PC est relié au commutateur \lstinline{"S1"} ; \lstinline{"S1"}
est relié à \lstinline{"S2"} et \lstinline{"S3"} ; \lstinline{"S2"} et \lstinline{"S3"} sont
aussi reliés entre eux (un cycle) ; le serveur est branché sur \lstinline{"S3"}. Le
dictionnaire a cinq clés et chaque câble est écrit des deux côtés :
% PYTHON-SOURCE: code/1NSI-RES-ME-003.py
\begin{python}
reseau = {
    "PC": ["S1"],
    "S1": ["PC", "S2", "S3"],
    "S2": ["S1", "S3"],
    "S3": ["S1", "S2", "Serveur"],
    "Serveur": ["S3"],
}
\end{python}
De \lstinline{"PC"} vers \lstinline{"Serveur"} : visités $\{\text{PC}\}$, on passe à
\lstinline{"S1"} ; visités $\{\text{PC}, \text{S1}\}$, on explore \lstinline{"S2"} ; de là
\lstinline{"S3"} n'est pas visité, on y va ; \lstinline{"Serveur"} est un voisin de
\lstinline{"S3"} : \lstinline{True}. Le cycle \lstinline{S1-S2-S3} n'a posé aucun problème,
parce que \lstinline{"S1"} était déjà dans les visités quand on l'a revu depuis \lstinline{"S3"}.

\textbf{Erreur fréquente.} Noter un câble d'un seul côté : avec
\lstinline|{"Box": ["PC", "TV"], "PC": ["Box"], "TV": []}|, la Box « voit » la TV mais la TV
ne voit personne. Ou croire qu'il faut interdire les cycles ou renoncer à la récursion : ce
qui protège de la boucle infinie, c'est le marquage des visités, pas la forme du réseau.

\end{fichemethode}
